\scanpage{72}

et en termes des $f\ldots$ la formule, est. identique mais d'allure un peu moins sympathique

\[
(11)\qquad f_{ll'}(u)=l'^{-1}\bigl(f_l(u)\bigr)\cdot f_{l'}(u)
\qquad\text{i.e. }f_u(ll')=l'^{-1}f_u(l)\cdot f_u(l')
\]

\deleted{La formule (10) peut s'exprimer en disant que pour $l$ fixé, $u\in\Gamma$ variable,}

\[
\deleted{\unsure{(17)|(12)}\qquad g_l:\Gamma\longrightarrow\widehat{\pi}_1(X,y)}
\]

\deleted{est un $1$-cocycle sur $\Gamma$, à valeurs dans $\widehat{\pi}_1(X,y)$, relativement à l'action naturelle de $\Gamma$ sur $\widehat{\pi}_1(X,y)$. Bien entendu, le $\widehat{\pi}_1(X,y)$-torseur à gauche, \inserted{\unsure{avec} op. de $\Gamma$,} correspondant à ce $1$-cocycle, n'est autre que $\operatorname{Hom}_{\widehat{\Pi}_1(X)}(x,y)$, sur lequel $\widehat{\pi}_1(X,y)=\operatorname{Hom}_{\widehat{\Pi}_1(X)}(y,y)$ opère par composition}

Pour un composé d'une \deleted{\unsure{nb} fini} suite finie composable de chemins

\diagram{\begin{tikzcd}[ampersand replacement=\&]
x_0\arrow[r,"l_1"]\&x_1\arrow[r,"l_2"]\&\cdots\&x_{n-1}\arrow[r,"l_n"]\&x_n
\end{tikzcd}}

la formule (10) prend la forme

\[
\unsure{(13)|(12)}\qquad
\begin{aligned}
g_{l_nl_{n-1}\mathinner{\ldotp\ldotp}l_1}(u)={}&g_{l_n}(u)\cdot l_n\bigl(g_{l_{n-1}}(u)\bigr)\cdot l_nl_{n-1}\bigl(g_{l_{n-2}}(u)\bigr)\mathinner{\ldotp\ldotp\ldotp\ldotp}\\
&\hfill \inserted{\cdot(l_n\mathinner{\ldotp\ldotp}l_2)g_{l_1}(u)}
\end{aligned}
\]

Formule la plus \deleted{évidente} \inserted{simple} possible en somme de celles qu'on pourrait imaginer qui \unsure{ait|ont} un sens. On va développer \deleted{\unsure{encore}} \inserted{un} contexte qui permette \deleted{\unsure{mieux|même}} d'écrire cette formule, ainsi que \deleted{(11)}, la formule analogue en $f$ et les formules

\margin{71}
